% ===========================================================================
%  Dynamic Hazard Rate BART -- software / technical documentation paper
%  Requires paper.bib in the same directory.
% ===========================================================================
\documentclass[11pt]{article}

\usepackage[utf8]{inputenc}
\usepackage[T1]{fontenc}
\usepackage{lmodern}
\usepackage[a4paper,margin=1in]{geometry}
\usepackage{amsmath,amssymb}
\usepackage{booktabs}
\usepackage{xcolor}
\usepackage{caption}
\usepackage{enumitem}
\usepackage{natbib}
\usepackage[colorlinks=true,linkcolor=teal,citecolor=teal,urlcolor=teal]{hyperref}

\title{\textbf{Dynamic Hazard Rate BART}\\[2pt]
\large An open-source platform for calibration-sensitive behavioral risk assessment}
\author{Ahmet Selim Y\i lmaz\\[2pt]
  \small Middle East Technical University, Ankara, T\"urkiye\\
  \small ORCID: \href{https://orcid.org/0009-0000-2866-5774}{0009-0000-2866-5774}}
\date{20 June 2026}

\begin{document}
\maketitle

\begin{center}
\small
\begin{tabular}{rl}
\textbf{Repository:} & \url{https://github.com/aslmylmz/metu-risk-persona}\\
\textbf{Archive (DOI):} & \href{https://doi.org/10.5281/zenodo.20592164}{10.5281/zenodo.20592164}\\
\textbf{Live platform:} & \url{https://bart.aselimyilmaz.com}\\
\textbf{License:} & MIT\\
\end{tabular}
\end{center}

\begin{abstract}
\noindent The Balloon Analogue Risk Task (BART) is a widely used behavioral
measure of risk-taking, but its standard parameterization places the
expected-value (EV) optimal stopping point near half of a balloon's maximum
capacity---so high that almost no participant reaches it, so the task cannot
separate calibrated risk-taking from gross exposure. \emph{Dynamic Hazard Rate BART}
replaces the classic uniform burst model with a sequence of independent
trials whose per-pump hazard rises linearly, $P(\text{burst at pump }k)=k/N$.
Under this model the EV-optimal stop is approximately $\sqrt{N}$, an attainable
target, restoring the task's ability to distinguish cautious, calibrated, and
risk-seeking play. The software bundles a React/Next.js game client, a Python scoring
engine that derives more than thirty psychometric metrics with a built-in
session-validation pipeline, and Monte Carlo tooling that verifies the analytic
optima. It is released under the MIT license and archived on Zenodo.
\end{abstract}

\section{Summary}

The Balloon Analogue Risk Task \citep{lejuez2002} is among the most widely used
behavioral measures of risk-taking. In its standard form, a balloon's burst point
is drawn once from a uniform distribution over its capacity $N$ (typically
$N=128$), which places the EV-optimal stopping point near $N/2 \approx 64$ pumps,
while participants stop systematically earlier. That gap is not a defect; it is a
valid index of risk aversion. The limitation that matters here is narrower:
because the optimum is rarely approached, the number of pumps and the money
earned are nearly collinear, so the task cannot separate calibrated risk-taking
from gross exposure \citep{frey2017}.

\emph{Dynamic Hazard Rate BART} is an open-source platform that restores the
measurement of calibration. It replaces the single uniform draw with a sequence
of independent trials whose per-pump hazard increases linearly,
$P(\text{burst at pump }k)=k/N$. Under this model the EV-optimal stop is
approximately $\sqrt{N}$---a reachable value---so the task can distinguish a
calibrated strategy from indiscriminate pumping and sort participants as cautious,
calibrated, or risk-seeking. The repository contains
three reusable parts: (i)~a React/Next.js game client that administers a
30-balloon, three-hazard session and logs pump-level telemetry; (ii)~a Python
scoring engine that derives more than thirty metrics (EV calibration, learning and
adaptation, post-explosion adjustment, response consistency, and a narrative
behavioral profile) with an explicit right-censoring correction and a
session-validation pipeline; and (iii)~Monte Carlo tooling that verifies the
analytically derived optima.

\section{Statement of need}

A persistent obstacle in the psychology and economics of risk is the
attitude--behavior gap: self-reported risk preferences correlate only weakly with
behavioral tasks, and behavioral tasks correlate weakly with one another
\citep{frey2017,pedroni2017}. A structural contributor to this gap is task design.
When the EV-optimal threshold is unreachable, the BART cannot reward
calibration---it can only index how far a participant pushed, not how well they
calibrated. Simulation work has
shown that the hazard structure of a BART systematically shapes the behavioral
profiles it can reconstruct \citep{diplinio2022}, yet most deployed
implementations retain the classic uniform model, and reusable, openly documented
scoring engines are scarce.

This platform addresses three needs for researchers who study risk behavior:

\begin{enumerate}[leftmargin=1.4em]
  \item \textbf{A calibration-sensitive task.} The linear-hazard model has a
  closed-form, psychologically attainable optimum ($s^*\approx\sqrt{N}$), enabling
  the task to distinguish cautious, calibrated, and risk-seeking
  participants \citep{pleskac2008,wallsten2005}.
  \item \textbf{A transparent, reusable scoring engine.} The engine computes
  EV-referenced calibration, learning, adaptation, and consistency metrics
  directly from raw event telemetry, with an explicit right-censoring
  (RNG-truncation) correction that uses collected (non-exploded) balloons for all
  behavioral-intention metrics.
  \item \textbf{Reproducibility and data quality.} A validation pipeline flags
  incomplete, too-fast, non-monotonic, automated, and key-repeat sessions, and a
  Monte Carlo script reproduces the analytic optima and the earnings distribution
  that anchors the engine's efficiency metric.
\end{enumerate}

The platform was built as the measurement instrument for a study at the Middle
East Technical University relating dynamic-hazard BART telemetry to the
domain-specific risk-taking (DOSPERT) scale \citep{weber2002,blais2006} in a
laboratory and online sample, with a pre-registered external replication against
the open Basel--Berlin Risk Study \citep{frey2017}.

\section{The dynamic-hazard model}

The redesign targets the hazard structure. Rather than treating risk as a bet on
a fixed future state, it models risk-taking as continuous adaptation to new
information: pursuing a larger reward means pumping further, which raises the
per-pump hazard the participant faces. For a balloon of capacity $N$, the
probability of surviving $s$ successive pumps is the product of the per-pump
survival probabilities,
%
\begin{equation}
P(\text{survive } s) = \prod_{k=1}^{s}\left(1 - \frac{k}{N}\right),
\end{equation}
%
and, with a reward of one unit per pump, the expected value of intending to stop
at $s$ is
%
\begin{equation}
\mathrm{EV}(s) = s \prod_{k=1}^{s}\left(1 - \frac{k}{N}\right).
\end{equation}
%
Taking the logarithm of the survival product and keeping the leading term gives
$\ln P(\text{survive } s) \approx -s^2/2N$, so survival is approximately Gaussian
in $s$ and $\mathrm{EV}(s) \approx s\,e^{-s^2/2N}$. Differentiating and setting the
result to zero yields the closed-form optimum
%
\begin{equation}
s^* = \sqrt{N}.
\end{equation}
%
The continuous approximation omits higher-order terms that make survival decay
slightly faster, so the realized integer optima are the floor of $\sqrt{N}$.
Table~\ref{tab:optima} lists the exact discrete optima for the three colors, which
a Monte Carlo simulation of 100{,}000 optimal sessions reproduces exactly.

\begin{table}[t]
\centering
\caption{The three balloon profiles and their exact discrete EV-optima.}
\label{tab:optima}
\begin{tabular}{lccccc}
\toprule
Color & Capacity $N$ & Risk tier & $s^*$ & $\mathrm{EV}(s^*)$ & $P(\text{survive } s^*)$\\
\midrule
Purple & 128 & Low    & 11 & 6.46 & 0.59\\
Teal   & 32  & Medium & 5  & 3.04 & 0.61\\
Orange & 8   & High   & 2  & 1.31 & 0.66\\
\bottomrule
\end{tabular}
\end{table}

Each banked pump pays \$0.25, so overshooting the optimum is penalized
endogenously through forfeited earnings rather than through an arbitrary
programmatic cap. A session presents 30 balloons---10 of each color in a shuffled
order---so the engine measures calibration \emph{to the risk level} rather than
gross risk-taking alone.

\section{Software architecture}

\paragraph{Game client.} A self-contained React component
(\texttt{games/bart/BartGame.tsx}) renders the 30 shuffled balloons, implements the
same $k/N$ hazard as the engine, records high-resolution monotonic timestamps via
\texttt{performance.now()}, and submits a typed session payload to a scoring
endpoint.

\paragraph{Scoring engine.} The Python engine (\texttt{scoring/bart.py}) segments
the event log into balloons, excludes operating-system key-repeat artifacts, and
computes EV-ratio calibration (reward-weighted across hazard levels), an explosion
penalty reported separately from calibration, three complementary learning
estimators, within- and between-balloon consistency, and a composite
adaptive-strategy score, before classifying each session into one of ten narrative
risk styles. The single most important methodological choice is a right-censoring
correction: because an exploded balloon records the random burst point rather than
the participant's intended stop, all behavioral-intention metrics are computed
from collected balloons only. Input and output are validated with pydantic models
\citep{pydantic}.

\paragraph{Verification and data tooling.} \texttt{scripts/monte\_carlo\_ev.py}
simulates 100{,}000 optimal sessions to confirm the integer optima and to
characterize the earnings distribution, and \texttt{scripts/generate\_synthetic.py}
produces synthetic datasets for testing and demonstration. The numerical core
relies on NumPy \citep{harris2020} and SciPy \citep{virtanen2020}.

\section{Quality control}

Remote behavioral data is vulnerable to bots, distraction, and superficial
engagement. The engine embeds a validation pipeline that marks a session invalid
when it has fewer than 15 of 30 balloons or contains out-of-order timestamps, and
that warns on incomplete sessions, color imbalance, suspiciously fast completion,
and near-zero pump variance. Separately, balloons whose pumps arrive at
operating-system key-repeat speed (median inter-pump interval below 80\,ms) are
detected and excluded from behavioral-intention metrics. These rules make the
engine's output auditable and reproducible across laboratory and online channels.

\section{Availability and documentation}

The source code is available on GitHub and archived on Zenodo
(\href{https://doi.org/10.5281/zenodo.20592164}{10.5281/zenodo.20592164}) under the
MIT license. Full documentation---installation, a quick-start scoring example, the
mathematical derivation, the complete metric definitions, and the validation
rules---is published on Read the Docs.

\section*{Acknowledgements}

This work was conducted at the Middle East Technical University under the
supervision of Assoc.\ Prof.\ G\"ul\c{s}ah Karakaya. The laboratory study was
reviewed and approved by the METU Human Subjects Ethics Committee (protocol
0176-ODT\"U\.{I}AEK-2026).

\bibliographystyle{plainnat}
\bibliography{paper}

\end{document}
